\documentclass[10pt,a4paper]{amsart}
\usepackage[margin=19mm]{geometry}
\usepackage{amsmath,amssymb,mathtools}
\usepackage[scr=rsfs,cal=euler]{mathalfa}
\usepackage{xfrac}
\usepackage[unicode,hidelinks,bookmarks=false]{hyperref}
\newcommand{\ie}{i.e.\ }
\newcommand{\cf}{cf.\ }
\newcommand{\resp}{respectively\ }
\newcommand{\Subsection}{\subsection}
\providecommand{\nobreakdash}{\nobreak\hbox{-}}
\setlength{\emergencystretch}{2em}
\multlinegap0pt
\usepackage{tikz}
\usetikzlibrary{arrows,positioning,shapes.geometric}
\newtheorem{theorem}{Theorem}
\newtheorem{lemma}{Lemma}
\newtheorem{proposition}{Proposition}
\newtheorem{corollary}{Corollary}
\newtheorem{remark}{Remark}
\newtheorem{definition}{Definition}
\begin{document}
\title[Global Non-Identifiability of Fubini-Study Geometry from Com]{Global Non-Identifiability of Fubini-Study Geometry from Complete One-Period Endpoint Data}
\author{Rashid Ahmad}
\date{}
\maketitle
\noindent\emph{Source and licence.} Global Non-Identifiability of Fubini-Study Geometry from Complete One-Period Endpoint Data. Version arXiv:2608.10051v1 (CC BY 4.0). This material is distributed under the Creative Commons Attribution 4.0 International licence, http://creativecommons.org/licenses/by/4.0/, as recorded in the arXiv OAI licence metadata for this exact version and shown on the abstract page https://arxiv.org/abs/2608.10051v1. Full rights notice: licenses/arxiv-2608.10051-CC-BY-4.0.txt.\par\smallskip
\noindent\textit{Editorial scope.} Counted unit: the original \texttt{theorem} environment \texttt{eq:G\_M} at source lines 434--459 together with its complete proof at lines 461--520; the delivered document keeps source lines 348--521 so that every referenced label and every citation resolves inside it. Native editable source is the paper's own concatenated TeX; the AMS wrapper is editorial formatting only. No publisher class, upstream script or unrelated artwork is redistributed.
\section{Observation Geometry and Exact Fibre Characterization}

The structure of the starting-time-indexed endpoint data follows directly from the Floquet relation.

\begin{proposition}[Starting-time-indexed observation]
For every $U\in\mathcal P_T^\infty$,
\begin{align}
E_{\rm ph}[U](\theta,[\tau])
&=
U_\theta(\tau)M_\theta U_\theta^\dagger(\tau),
\label{eq:orbit_observation}
\end{align}
where $M_\theta=U_\theta(T)$.
\end{proposition}

\begin{proof}
Using the composition law and the Floquet relation,
\begin{align}
U_\theta(\tau+T,\tau)
&=
U_\theta(\tau+T,0)
U_\theta^\dagger(\tau,0)
\nonumber\\
&=
U_\theta(\tau)M_\theta
U_\theta^\dagger(\tau).
\end{align}
\end{proof}

Define the observed conjugation-orbit path
\begin{align}
\Gamma_U(\theta,[t])
:=
\operatorname{Ad}_{U_\theta(t)}M_\theta
=
U_\theta(t)M_\theta U_\theta^\dagger(t).
\end{align}
Equation~\eqref{eq:orbit_observation} shows that
\begin{align}
E_{\rm ph}[U]=\Gamma_U.
\end{align}
Thus, for each $\theta$, the complete starting-time-indexed endpoint data determine the path of $M_\theta$ through its conjugacy orbit
\begin{align}
\mathcal O_{M_\theta}
:=
\{VM_\theta V^\dagger:V\in U(d)\},
\end{align}
but do not, in general, determine the unitary $U_\theta(t)$ that generates this path.

The resulting ambiguity is characterized by the elementary identity
\begin{align}
\operatorname{Ad}_{U_1(t)}M
=
\operatorname{Ad}_{U_2(t)}M
\quad\Longleftrightarrow\quad
[U_2^\dagger(t)U_1(t),M]=0.
\label{eq:orbit_fibre_identity}
\end{align}
Hence the relevant fibres are generated by right multiplication with paths taking values in the centralizer of the monodromy.

An infinitesimal form of the same observation follows by differentiating the orbit path:
\begin{align}
\dot\Gamma(t)
&=
\dot U(t)MU^\dagger(t)
+
UM\dot U^\dagger(t)
\nonumber\\
&=
-i[H(t),\Gamma(t)].
\label{eq:Gamma_dot}
\end{align}
Consequently, the observed tangent vector determines only the commutator component of the Hamiltonian. Its kernel is the Lie-algebra centralizer
\begin{align}
\ker\operatorname{ad}_{\Gamma(t)}
=
\mathfrak z(\Gamma(t)).
\end{align}
Since
\begin{align}
\mathfrak z(\Gamma(t))
=
U(t)\mathfrak z(M)U^\dagger(t),
\end{align}
the infinitesimal ambiguity is precisely the tangent-space manifestation of the global centralizer fibre. The infinitesimal relation is useful for interpretation but is not required for the global fibre characterization.


\begin{theorem}[Fixed-monodromy fibre characterization]
Fix $M\in C^\infty(I,U(d))$. On $\mathcal P_M^\infty$, define
\begin{align}
\mathcal G_M^0
:=
\Bigl\{
V\in C^\infty(I\times\mathbb R,U(d)):
&\;V_\theta(0)=\mathbb I,
\nonumber\\
&\;V_\theta(t+T)=V_\theta(t),
\nonumber\\
&\;[V_\theta(t),M_\theta]=0
\quad\forall(\theta,t)
\Bigr\}.
\label{eq:G_M}
\end{align}
Then, for $U_1,U_2\in\mathcal P_M^\infty$,
\begin{align}
E_{\rm ph}[U_1]=E_{\rm ph}[U_2]
\quad\Longleftrightarrow\quad
U_1=U_2V
\quad\text{for some }V\in\mathcal G_M^0.
\label{eq:fibre_characterization}
\end{align}
Thus every fibre of $E_{\rm ph}$ on $\mathcal P_M^\infty$ is a right $\mathcal G_M^0$-orbit.
\end{theorem}

\begin{proof}
Suppose first that
$E_{\rm ph}[U_1]=E_{\rm ph}[U_2]$. By Eq.~\eqref{eq:orbit_observation},
\begin{align}
U_1(t)MU_1^\dagger(t)
=
U_2(t)MU_2^\dagger(t).
\end{align}
Define
\begin{align}
V(t):=U_2^\dagger(t)U_1(t).
\end{align}
Then
\begin{align}
V(t)M
=
MV(t),
\end{align}
so $V(t)$ takes values in $Z(M)$. Since both propagators are normalized at $t=0$,
\begin{align}
V(0)=\mathbb I.
\end{align}
Furthermore,
\begin{align}
V(t+T)
&=
U_2^\dagger(t+T)U_1(t+T)
\nonumber\\
&=
(U_2(t)M)^\dagger(U_1(t)M)
\nonumber\\
&=
M^\dagger V(t)M
=
V(t),
\end{align}
where the last equality follows from $[V(t),M]=0$. Hence $V\in\mathcal G_M^0$ and $U_1=U_2V$.

Conversely, suppose $U_1=U_2V$ with $V\in\mathcal G_M^0$. Then
\begin{align}
U_1(t)MU_1^\dagger(t)
&=
U_2(t)V(t)MV^\dagger(t)U_2^\dagger(t)
\nonumber\\
&=
U_2(t)MU_2^\dagger(t),
\end{align}
so $E_{\rm ph}[U_1]=E_{\rm ph}[U_2]$. Moreover,
\begin{align}
U_1(t+T)
&=
U_2(t)MV(t)
\nonumber\\
&=
U_2(t)V(t)M
=
U_1(t)M,
\end{align}
and therefore $U_1\in\mathcal P_M^\infty$.
\end{proof}


\end{document}
